% Generated draft: requires mathematical and source review.
\documentclass[11pt]{article}
\usepackage[margin=25mm]{geometry}
\usepackage{fontspec}
\setmainfont{Latin Modern Roman}
\usepackage{amsmath,amssymb,mathtools,unicode-math}
\setmathfont{Latin Modern Math}
\usepackage{enumitem,needspace,xurl,hyperref}
\setlength{\parindent}{0pt}
\setlength{\parskip}{5pt}
\setlength{\emergencystretch}{3em}
\newcommand{\N}{\mathbb{N}}
\newcommand{\Z}{\mathbb{Z}}
\newcommand{\Q}{\mathbb{Q}}
\newcommand{\R}{\mathbb{R}}
\newcommand{\C}{\mathbb{C}}
\newcommand{\F}{\mathbb{F}}
\newcommand{\A}{\mathbb{A}}
\newcommand{\PP}{\mathbb{P}}
\newcommand{\E}{\mathbb{E}}
\newcommand{\eps}{\varepsilon}
\newcommand{\dd}{\,\mathrm{d}}
\begin{document}
\section*{UnsolvedMath ID: 701074}
\textbf{Source problem code:} catalogue entry
\subsection{Definitions and mathematical statement}
For a group \(P\), define its lower central series by
\[
\gamma_1(P)=P,\qquad \gamma_{i+1}(P)=[\gamma_i(P),P]\quad(i\geq1),
\]
where \([H,P]\) is the subgroup generated by commutators \(h^{-1}x^{-1}hx\) with \(h\in H\) and \(x\in P\). Define its upper central series by
\[
\zeta_0(P)=1,\qquad
\zeta_{j+1}(P)/\zeta_j(P)=Z(P/\zeta_j(P))\quad(j\geq0),
\]
where \(Z(H)\) denotes the center of a group \(H\). The nilpotency class \(c\) is the least nonnegative integer for which \(\gamma_{c+1}(P)=1\).

The source asks whether a finite \(p\)-group \(P\) of ``abundance'' \(a\) and nilpotency class \(c\) admits an integer bound \(t=t(a)\) such that
\[
\gamma_i(G)=\zeta_{c-i+1}(G)\qquad(i\geq t).
\]
Here the source does not define abundance, introduces \(G\) without identifying it with \(P\), and gives no upper bound on \(i\), although the upper central series index becomes negative when \(i>c+1\).

\emph{Conditional reconstruction requiring review.} If \(G\) means \(P\), and the intended range consists of the indices for which both displayed series terms are defined, the question becomes: for every admissible abundance value \(a\), does there exist an integer \(t(a)\geq1\), depending only on \(a\), such that for every prime \(p\) and every finite \(p\)-group \(P\) of abundance \(a\) and nilpotency class \(c\),
\[
\gamma_i(P)=\zeta_{c-i+1}(P)
\qquad\text{for every integer }i\text{ with }t(a)\leq i\leq c+1?
\]
The endpoint \(i=c+1\) gives the automatic equality \(1=1\). The definition and admissible values of abundance remain unspecified, so this reconstruction is not a fully determined mathematical statement.

The supplied source states that the assertion holds for \(a=0\) with \(t=1\), because \(P\) then has maximal class, and that for \(a=1\) one can take \(t=3\). These are preserved as source assertions, without independent verification.
\subsection{Short English statement}
The question asks whether abundance alone bounds how far down the lower central series a finite \(p\)-group must be before it agrees with the upper central series read in reverse. Precision requires a definition of abundance, clarification that \(G=P\), and an intended index range. The supplied registry labels the problem open; that label is not an independently verified status assessment.
\subsection{Sources}
\begin{itemize}
\item Local UnsolvedMath registry, permanent ID \(701074\), exact source problem code \texttt{catalogue entry}, title \emph{Kourovka Notebook Problem 14.90}; supplied status: \texttt{open}. Upstream revision: \texttt{d371ee917ac9647fbb06cd6fc529deec6f3dda82}.
\item \emph{The Kourovka Notebook (Unsolved Problems in Group Theory)}, problem 14.90, attributed to G. Fern\'andez--Alcober, proposed in 1999; the supplied background identifies the 14th issue (1999) and the 21st issue edition (2026), arXiv:\texttt{1401.0300v48}. The collection description names E. I. Khukhro and V. D. Mazurov as editors.
\item The statements concerning abundance \(0\) and \(1\) are supplied with the problem text. No separate proofs or supporting references for these statements are supplied. The cited notebook has not been independently consulted for this draft.
\end{itemize}
\end{document}
